\chapter{FPGA տվյալների գեներացիա և մշակում}

Սկզբնական դիպլոմային տարբերակում իրական սենսորները պարտադիր չեն։ Միկրոէլեկտրոնիկայի
ուսանողը պետք է իրականացնի \textbf{կոնֆիգուրելի ազդանշան-գեներատոր / սենսոր-էմուլյատոր}
RTL մոդուլներ, որոնք FPGA-ի ներսում մոդելավորում են սենսորի իրատեսական վարքը։
Հնարավոր սենսորների տիպերը՝ ջերմաստիճան (Temperature), ճնշում (Pressure),
վիբրացիա (Vibration), ընդհանրական անալոգային սենսոր և այլ կոնֆիգուրելի մոդելներ։

Յուրաքանչյուր սենսորի համար առանձին ապարատ իրականացնելու փոխարեն նախատեսվում են
\textbf{վերաօգտագործելի ազդանշան-գեներացման պրիմիտիվներ}.

\begin{itemize}[leftmargin=1.4em]
  \item \textbf{Ramp Generator} — աստիճանաբար աճող/նվազող արժեքներ (ջերմաստիճանի,
        ճնշման, դանդաղ փոփոխվող մեծությունների համար)։
  \item \textbf{Periodic Signal Generator} — պարբերական ազդանշաններ (վիբրացիա և
        այլ կրկնվող երևույթներ)։
  \item \textbf{Pseudo-Random / Noise Generator} — LFSR-ի հիման վրա, կոնֆիգուրելի
        աղմուկ ավելացնելու համար։
  \item \textbf{Anomaly / Fault Generator} — ոչ նորմալ վարքի ներարկում (օր.՝
        գերտաքացման կամ բարձր վիբրացիայի իրադարձություն), ինչը տվյալները դարձնում է
        օգտակար անոմալիաների հայտնաբերման փորձերի համար։
\end{itemize}

\adddiagram[fig:rtl]{rtl-pipeline}{0.34}{FPGA-ի RTL մշակման հնարավոր փոխակրիչը (pipeline).}

\textbf{Signal Generator}-ի կոնֆիգուրելի պարամետրերն են՝ ազդանշանի տիպ, ամպլիտուդ,
հաճախություն, բազային արժեք, նմուշառման արագություն (sample rate), աղմուկի մակարդակ,
անոմալիայի միացում և մեծություն։ \textbf{Sensor Emulator}-ը ընդհանրական ազդանշանը
վերածում է սենսորանման արժեքների՝ scaling, offset, saturation, աղմուկի ու անոմալիայի
ներարկման միջոցով։

FPGA-ն պետք է կատարի հիմնական իրական ժամանակի մշակում (\textbf{Processing IP}-ներ).
հնարավոր մոդուլներ՝ Moving Average, Min/Max Detector, Threshold Detector, Event
Detector, FIR Filter, Statistical accumulator։ Բակալավրիատի ծավալի համար ողջամիտ
MVP-ն պարունակում է մոտավորապես \textbf{3 մշակող մոդուլ}։
